% !TeX root = ../深度学习与大语言模型技术原理.tex
% 由主文档通过 \input 载入；请编译主文档。
\setcounter{chapter}{21} % 保留原书章号。
\chapter{自回归推理与键值缓存}
\label{ch:inference}

\section{预填充与逐词元解码}

推理先处理输入提示，称为预填充（prefill），得到各层表示及首个待生成词元的分布；随后逐次追加词元，称为解码（decode）。预填充可并行处理提示位置，解码的每一步依赖上一步生成结果。贪心解码选择最高概率词元，随机解码按指定分布采样；贪心局部最优不保证整个序列概率最大，也不保证任务正确。

预填充通常包含较大的矩阵乘法，解码在小批量时需要为少量新词元反复读取大量权重与历史缓存，因此常呈现计算与带宽瓶颈的差异。但批量足够大、上下文很长或采用不同架构时，瓶颈可能变化，应根据实测算术强度判断，而非把“prefill 计算受限、decode 带宽受限”当作绝对定律。

首词元延迟（time to first token, TTFT）包含排队、调度、预填充和首个输出准备。词元间延迟（inter-token latency, ITL）描述相邻输出间隔，平均每输出词元时间（time per output token, TPOT）常按首词元后生成耗时平均。吞吐量可按输入词元、输出词元或请求数定义；必须写清分母和统计窗口，避免把不同口径的数字直接比较。

\section{缓存为什么有效}

对确定性、因果的解码器，在参数、历史词元、位置编号和归一化规则不变时，追加未来词元不会改变过去位置的隐藏状态。每层过去位置的键和值因而可以保存，新词元只计算自己的查询、键和值。查询用于本次访问历史，通常不需要作为后续层的长期注意力缓存。

第 $t$ 步的新查询访问缓存 $K_{1:t},V_{1:t}$，其注意力算术成本仍随历史长度增长。KV Cache 避免历史位置的重复投影与 FFN，也避免每步重算完整历史的注意力；它没有让每个新词元的注意力成为常数成本。若从空前缀生成 $T$ 个词元，固定维度下缓存注意力总成本为 $O(T^2d)$，而每步重算全序列的朴素注意力总成本为 $O(T^3d)$。

带 RoPE 时通常缓存已旋转的键，追加位置要使用正确绝对偏移。若中途改变位置缩放或模型权重，旧缓存可能失效。滑动窗口、截断、分支生成和多模态位置编码也要求专门管理，不能仅在数组末尾追加便假定正确。

\section{缓存存储与批量容量}

在每层都采用相同头配置、键值维度相同、每元素占 $b$ 字节且批内长度一致时，缓存数据量为
\begin{equation}
M_{\mathrm{KV}}=2BLnh_{kv}d_hb.
\label{eq:kv-memory}
\end{equation}
系数2来自键和值。变长批量可把 $Bn$ 替换为有效长度之和，但实际还要加块填充、索引、尺度元数据和临时空间。张量并行可能分片缓存，也可能因头数与分片数关系发生复制，应按实现核算。

\begin{example}
设 $B=1,L=32,n=8192,h_{kv}=8,d_h=128$，键值均为2字节元素，求缓存容量；若改为32个键值头，容量如何变化？
\end{example}
\begin{solve}
代入式~\eqref{eq:kv-memory} 得 $1073741824$ 字节，即1 GiB。键值头从8增至32时，数据量增为4 GiB。该结果尚未包括模型权重和其他缓冲，展示了 GQA 对长上下文容量的直接影响。
\end{solve}

\section{增量解码的正确性检查}

给定相同模型、词元和位置，在关闭 dropout 后，比较一次全序列前向与逐步缓存前向在对应位置的 logits，应在合理数值容差内一致。不同内核与归约顺序可能造成小误差，随机采样又可能把很小分布差异转成不同文本，所以检查 logits 比只比较最终句子更有诊断性。

变长批量需要区分物理存储位置、逻辑位置与填充位置。缓存复用时还要检查模型版本、适配器、前缀词元和位置配置。若使用共享缓存，某个请求修改尾块不能破坏其他请求，后文的写时复制正是为此服务。

